% authority: science_draft
\documentclass[11pt]{article}
\usepackage[T1]{fontenc}
\usepackage{lmodern}
\usepackage[margin=1in]{geometry}
\usepackage{amsmath,amssymb,amsthm,mathtools}
\usepackage{booktabs,longtable,array,enumitem}
\newcommand{\Qadm}{\mathcal Q_{\mathrm{adm}}}
\newcommand{\Red}{\mathsf{Red}_{\mathrm{src}}}
\newcommand{\PhiSrc}{\Phi_{\mathrm{src}}}
\newcommand{\Lam}{\Lambda_+}
\newcommand{\Qbar}{\overline{\mathcal Q}_{\mathrm{adm}}}
\newtheorem{definition}{Definition}
\newtheorem{theorem}{Theorem}
\newtheorem{proposition}{Proposition}
\newtheorem{boundary}{Boundary}
\title{V21 P5-T02: Proposal-Only Semantics for \(\Phi_{\mathrm{src}}\)}
\author{Ontology Formalizer parent--child synthesis}
\date{2026-07-25}

\begin{document}
\maketitle

\begin{center}
\fbox{\parbox{0.93\linewidth}{\textbf{Status and authority.}
This source is \texttt{draft/control}, \texttt{proposal-only}, and
\texttt{source-extension data}, with disposition
\texttt{blocked\_adoption\_open\_continuation}. It tests one source-side
candidate law; it does not adopt, reject, or modify canonical ontology. The
parameter below is not established as physical time, a clock, a gauge
parameter, or physical evolution. No source dynamics, reconstruction, metric,
matter coupling, Einstein equation, benchmark promotion, proof publication,
push, or completed derivation follows.}}
\end{center}

\section{Source basis and exact question}

The P5-T01 compact source-theory object contains admissible source
configurations \(\Qadm\), declared source redundancy \(\Red\), and a tagged
open semantics slot \(S_\Phi\). It deliberately provides no value for that
slot. P5-T02 asks which mathematical kind is sufficiently definite to support
later source-dynamics work without importing a target atlas, target metric,
benchmark behavior, or physical interpretation.

This packet tests one conservative answer: \(\PhiSrc\) is a continuous unital
forward action of an ordered additive topological monoid. The word
``semiflow'' below names this mathematical structure only. It is not a claim
that the configurations are physical states or that the parameter is time.

\section{Selected candidate semantic type}

\begin{definition}[Forward parameter monoid]
Let
\[
  \Lam=(\mathbb R_{\geq 0},+,0,\leq)
\]
with its Euclidean topology. This is a commutative, ordered, topological
monoid. It contains no negative parameters.
\end{definition}

\begin{definition}[P5-T02 source semiflow candidate]
Give \(\Qadm\) a declared source-only topology \(\mathcal T_{\mathrm{adm}}\).
A P5-T02 candidate is a map
\[
  \PhiSrc:\Lam\times\Qadm\longrightarrow\Qadm,
  \qquad (\lambda,q)\longmapsto\Phi_\lambda(q),
\]
with the following laws:
\begin{enumerate}[label=(S\arabic*)]
  \item \textbf{Identity:} \(\Phi_0(q)=q\).
  \item \textbf{Forward composition:}
  \[
    \Phi_{\lambda+\mu}(q)
    =\Phi_\lambda(\Phi_\mu(q)).
  \]
  Thus applying \(\mu\) and then \(\lambda\) agrees with the summed
  parameter. This is the conventional left-action notation. Commutativity of
  addition also implies
  \(\Phi_\lambda\circ\Phi_\mu=\Phi_\mu\circ\Phi_\lambda\), but the displayed
  convention is fixed to prevent later notational drift.
  \item \textbf{Admissibility:}
  \(\Phi_\lambda(q)\in\Qadm\) for every
  \((\lambda,q)\in\Lam\times\Qadm\).
  \item \textbf{Redundancy compatibility:}
  \[
    q\sim_{\mathrm{red}}q'
    \Longrightarrow
    \Phi_\lambda(q)\sim_{\mathrm{red}}\Phi_\lambda(q')
  \]
  for every \(\lambda\geq0\).
  \item \textbf{Joint continuity:} \(\PhiSrc\) is continuous as a map from
  \(\Lam\times(\Qadm,\mathcal T_{\mathrm{adm}})\) to
  \((\Qadm,\mathcal T_{\mathrm{adm}})\).
\end{enumerate}
\end{definition}

The candidate fills only the mathematical value of \(S_\Phi\), and only at
proposal scope. The P5-T01 slots \(S_{\mathrm{dyn}}\) and
\(S_{\mathrm{rec}}\) remain tagged open. The selected kind is a monoid action,
not a group action: the domain contains no negative parameter and the laws do
not require any \(\Phi_\lambda\) to be invertible.

\begin{boundary}[Source topology and no target import]
The topology \(\mathcal T_{\mathrm{adm}}\), admissibility rule, parameter
order, and redundancy relation must be supplied from source-side definitions.
They may not be chosen by matching a target atlas, target metric, exact-GR
solution, benchmark success, or candidate reconstruction.
\end{boundary}

\section{Basic law consequences}

\begin{proposition}[Finite composition]
For every \(n\geq1\), \(\lambda_1,\ldots,\lambda_n\in\Lam\), and
\(q\in\Qadm\),
\[
  \Phi_{\lambda_n}\circ\cdots\circ\Phi_{\lambda_1}(q)
  =
  \Phi_{\lambda_1+\cdots+\lambda_n}(q).
\]
\end{proposition}

\begin{proof}
The case \(n=1\) is immediate. If the formula holds for \(n\), axiom (S2)
gives
\[
 \Phi_{\lambda_{n+1}}
 \bigl(\Phi_{\lambda_1+\cdots+\lambda_n}(q)\bigr)
 =
 \Phi_{\lambda_1+\cdots+\lambda_{n+1}}(q).
\]
This proves the result by induction.
\end{proof}

\begin{theorem}[Redundancy-quotient descent]
Let \(\pi:\Qadm\to\Qbar=\Qadm/{\sim_{\mathrm{red}}}\) be the orbit projection.
Under (S4), the formula
\[
  \overline\Phi_\lambda([q])
  =
  [\Phi_\lambda(q)]
\]
defines a unique action of \(\Lam\) on the quotient set \(\Qbar\). It obeys
identity and forward composition. If \(\Qbar\) has the quotient topology and
\(\mathrm{id}_{\Lam}\times\pi\) is a quotient map---for example, under a
separately checked sufficient condition such as an open orbit
projection---then the induced map
\(\overline\Phi:\Lam\times\Qbar\to\Qbar\) is continuous.
\end{theorem}

\begin{proof}
If \([q]=[q']\), axiom (S4) gives
\([\Phi_\lambda(q)]=[\Phi_\lambda(q')]\), so the formula is well defined.
Surjectivity of \(\pi\) makes it unique. Applying (S1) and (S2) before taking
equivalence classes proves identity and composition. For the topological
claim,
\[
 \overline\Phi\circ(\mathrm{id}_{\Lam}\times\pi)
 =
 \pi\circ\PhiSrc .
\]
The right side is continuous. The stated quotient-map hypothesis therefore
gives continuity of \(\overline\Phi\). Without that extra topological
hypothesis only the algebraic descent is claimed.
\end{proof}

This theorem is relative to the declared redundancy relation. It does not
prove that \(\Red\) is a physical gauge symmetry or that the quotient is
physical state space. Because the action is on \(\Qadm\), it also supplies no
action on the source arena, no bundle lift, and no trajectory of source
points.

\section{Reparameterization and the clock nonconclusion}

\begin{theorem}[Continuous additive reparameterization classification]
Let \(r:\mathbb R_{\geq0}\to\mathbb R_{\geq0}\) be a continuous,
order-preserving monoid endomorphism:
\[
 r(0)=0,\qquad r(\lambda+\mu)=r(\lambda)+r(\mu).
\]
Then there is \(c\geq0\) such that \(r(\lambda)=c\lambda\) for all
\(\lambda\geq0\). It is an order-preserving monoid automorphism exactly when
\(c>0\).
\end{theorem}

\begin{proof}
Set \(c=r(1)\). Additivity gives \(r(n)=nc\) for
\(n\in\mathbb N_0\), and \(n\,r(m/n)=r(m)=mc\), so
\(r(m/n)=c(m/n)\) for nonnegative rationals. Approximate any
\(\lambda\geq0\) by nonnegative rationals and use continuity to obtain
\(r(\lambda)=c\lambda\). Nonnegativity gives \(c\geq0\). For \(c>0\), the
inverse is \(\lambda\mapsto\lambda/c\); \(c=0\) collapses every parameter and
is not an automorphism.
\end{proof}

The reparameterized family
\(\Phi^{(c)}_\lambda=\Phi_{c\lambda}\) therefore obeys the same monoid laws
for \(c\geq0\), with faithful scale change for \(c>0\). The selected semantics
does not determine \(c\), a unit, a rate normalization, or any relation to a
physical clock. This residual scaling freedom is a mathematical reason not to
read the parameter as measured time.

\section{Invertibility and generator boundaries}

The mainline candidate requires no inverse. A particular
\(\Phi_\lambda:\Qadm\to\Qadm\) may happen to be bijective, but its inverse
need not occur at a parameter in \(\Lam\). A global group action is available
only on a separately proved invariant domain \(D\subseteq\Qadm\) when the
forward family extends consistently to
\(\widetilde\Phi:\mathbb R\times D\to D\). Neither bijectivity nor such an
extension is supplied here.

Continuity alone also gives no infinitesimal generator. A generator such as
\[
 X(q)=\left.\frac{d}{d\lambda}\Phi_\lambda(q)\right|_{\lambda=0^+}
\]
requires a separately declared differentiable or functional-analytic
structure, existence of the derivative, and domain control. Those are later
burdens. No equation of motion follows from the present action laws.

\section{Witnesses and specializations}

\subsection{Continuous absorbing-translation witness}

Take \(\Qadm=\mathbb R_{\geq0}\) with its ordinary topology, identity-only
declared redundancy, and
\[
  \Phi_\lambda(x)=\max\{x-\lambda,0\}.
\]
Then \(\Phi_0=\mathrm{id}\),
\[
 \max\{\max\{x-\lambda,0\}-\mu,0\}
 =\max\{x-(\lambda+\mu),0\},
\]
and the map is jointly continuous. For every \(\lambda>0\), distinct points
in \([0,\lambda]\) map to \(0\), so the action is not invertible. This proves
that the selected axioms are consistent and do not secretly imply a group.
It is a mathematical witness, not an empirical source model.

\subsection{Discrete specialization}

For the discrete parameter monoid \((\mathbb N_0,+,0)\), take
\(\Qadm=\{a,b,c\}\) and
\[
 F(a)=b,\qquad F(b)=b,\qquad F(c)=c,\qquad \Phi_n=F^n.
\]
Let declared redundancy identify \(a\sim_{\mathrm{red}}b\) and leave \(c\)
alone. The action is unital, associative, noninvertible, and compatible with
the redundancy, so it descends to the two-class quotient. This is a discrete
specialization of the algebraic law, not a continuous instance of the
\(\mathbb R_{\geq0}\) mainline and not a physical discretization claim.

\section{Required exceptional-case classification}

\begin{longtable}{p{0.18\linewidth}p{0.40\linewidth}p{0.32\linewidth}}
\toprule
Case & Mathematical disposition & Forbidden overread \\
\midrule
\endhead
Continuous parameter
& Selected mainline: \(\mathbb R_{\geq0}\) with joint continuity.
& Continuity does not supply a clock, rate, generator, or physical evolution.\\
Discrete parameter
& Admissible specialization using \(\mathbb N_0\); not equivalent to the
continuous mainline without an embedding or extension theorem.
& A step label is not a physical time interval.\\
Periodic orbit
& Allowed when \(\Phi_p(q)=q\) for some \(p>0\); a circle rotation gives a
simple mathematical example.
& The return parameter is not an observed physical period.\\
Bi-infinite orbit
& Not part of the forward mainline; it requires a separately proved
\(\mathbb R\)-action extension on an invariant subdomain.
& Backward physical evolution or reversibility is not inferred.\\
Dense orbit
& Compatible with a continuous action, for example on a torus with
incommensurate rotation rates.
& Density does not fix physical topology, ergodicity, or a clock.\\
Branching
& Excluded by the selected single-valued map. It requires a relation,
multiflow, or categorical extension and a new law packet.
& Exclusion from this candidate is not a global impossibility theorem.\\
Stochastic
& Excluded by the deterministic map. It requires transition kernels and a
Chapman--Kolmogorov or equivalent composition law.
& No probability, noise, measurement, or stochastic physics is inferred.\\
\bottomrule
\end{longtable}

\section{Fail-closed conditions}

The candidate is rejected as malformed for this packet if any of the
following occurs:
\begin{enumerate}
  \item target atlas, target metric, benchmark success, or candidate
  reconstruction defines a source component;
  \item identity, composition, admissibility, or redundancy compatibility
  fails;
  \item continuity is claimed without a declared source topology;
  \item quotient continuity is claimed without the required quotient-map
  hypothesis;
  \item a global inverse, group extension, or generator is asserted without
  its additional hypotheses;
  \item \(\lambda\) is identified with physical time, clock scale, gauge,
  physical dynamics, or a reconstruction parameter by the action axioms
  alone.
\end{enumerate}
Failure is scoped to this candidate or claim. It does not establish that a
different source-side extension is impossible.

\section{Burden disposition and exact nonconclusions}

At draft/control source-only scope, P5-T02 removes one kind of ambiguity:
it gives \(S_\Phi\) a precise candidate type and law set and proves two basic
structural results. The action remains proposal-only source-extension data.
Current adoption is blocked while same-milestone continuation remains open.
P5-T03 must supply a source-dynamics candidate compatible with this structure
or explain precisely why a different \(\PhiSrc\) law is required.

Nothing here establishes:
\begin{itemize}
  \item a canonical-ontology candidate, adopted or rejected source law, or
  physical source manifold \(M_{\mathrm{src}}\);
  \item physical time, clock normalization, gauge symmetry, reversibility,
  an infinitesimal generator, or an equation of motion;
  \item a source-to-target reconstruction, effective metric
  \(g_{\mathrm{eff}}\), causal structure, or target diffeomorphism symmetry;
  \item matter semantics, universal coupling, stress energy, or backreaction;
  \item Einstein-equation recovery or exact-GR benchmark promotion;
  \item proof publication, push, or completed-derivation authority.
\end{itemize}

\end{document}
